\documentclass[12pt,a4paper]{article}
\usepackage[utf8]{inputenc}
\usepackage[T1]{fontenc}
\usepackage[margin=1in]{geometry}
\usepackage{booktabs}
\usepackage{array}
\usepackage{siunitx}
\usepackage{multirow}
\usepackage{tabularx}
\usepackage{graphicx}
\usepackage{amsmath}
\usepackage{caption}
\usepackage{hyperref}

\title{Weighted A* Search for the N-Puzzle Problem}
\author{Sabbir Ahmmed \\ 2205040 \\ CSE 318 -- Offline 1} 
\date{}

\begin{document}
\maketitle

\section{Introduction}

The N-puzzle is a classic problem in artificial intelligence, where tiles must be
slid into their correct positions using the blank space. The A* search algorithm
is widely used to find optimal solutions to such problems. A* evaluates states
using the cost function
\[
f(n) = g(n) + h(n)
\]
where $g(n)$ is the cost to reach node $n$ from the start, and $h(n)$ is a
heuristic estimate of the remaining cost to the goal. The algorithm guarantees
an optimal solution when $h(n)$ is admissible (never overestimates) and
consistent.

\subsection*{Weighted A*}

Weighted A* modifies the evaluation function to
\[
f(n) = g(n) + W \cdot h(n)
\]
with $W \ge 1.0$. Higher values of $W$ place more emphasis on the heuristic,
making the search more greedy and faster, at the cost of optimality. Lower
values (close to $1.0$) behave like classic A* and produce optimal paths but
explore more nodes.

\section{Heuristic Design}

We combine two heuristics to create an informed and admissible estimate:

\begin{itemize}
  \item \textbf{Manhattan Distance}: For each tile, the sum of the absolute
  differences between its current and goal coordinates. This is admissible
  because each tile must move at least that many steps.

  \item \textbf{Linear Conflict}: When two tiles are in their correct row (or
  column) but in the wrong order relative to each other, at least one must move
  out and back, adding a penalty of 2 moves per conflict. This improves upon
  Manhattan distance while remaining admissible.
\end{itemize}

The combined heuristic is:
\[
h(n) = \text{Manhattan}(n) + \text{LinearConflict}(n)
\]

This heuristic is both admissible and consistent, ensuring that A* finds the
optimal solution when $W = 1.0$.

\section{Experimental Setup}

We generated 10 distinct solvable $4 \times 4$ puzzle configurations. For each
configuration, we ran Weighted A* with $W \in \{1.0, 1.2, 2.0, 5.0\}$ and
recorded the solution cost (number of moves) and the number of nodes explored.

\section{Results}

\subsection*{Summary Table}

\begin{table}[ht!]
\centering
\caption{Average metrics across 10 test cases for each weight value}
\label{tab:summary}
\begin{tabular}{lSSS}
\toprule
{Weight} & {Avg. Nodes Explored} & {Avg. Cost} & {Avg. Cost Ratio} \\
\midrule
1.0 & 61344.6 & 37.8 & 1.000 \\
1.2 & 12379.0 & 38.6 & 1.021 \\
2.0 & 1646.7 & 46.6 & 1.233 \\
5.0 & 422.9 & 60.6 & 1.603 \\
\bottomrule
\end{tabular}
\end{table}

\subsection*{Detailed Results}

\begin{table}[ht!]
\centering
\caption{Solution cost (number of moves) for each test case across different weights}
\label{tab:cost}
\begin{tabular}{c|SSSS}
\toprule
{Test} & \multicolumn{4}{c}{Weight $W$} \\
\cline{2-5}
{Case} & {1.0} & {1.2} & {2.0} & {5.0} \\
\midrule
0 & 46 & 46 & 60 & 62 \\
1 & 40 & 42 & 54 & 72 \\
2 & 44 & 44 & 56 & 70 \\
3 & 42 & 42 & 50 & 72 \\
4 & 28 & 28 & 28 & 56 \\
5 & 34 & 34 & 40 & 44 \\
6 & 44 & 46 & 52 & 74 \\
7 & 42 & 42 & 48 & 60 \\
8 & 28 & 30 & 36 & 52 \\
9 & 30 & 32 & 42 & 44 \\
\midrule
Avg & 37.8 & 38.6 & 46.6 & 60.6 \\
\bottomrule
\end{tabular}
\end{table}

\begin{table}[ht!]
\centering
\caption{Nodes explored for each test case across different weights}
\label{tab:nodes}
\begin{tabular}{c|SSSS}
\toprule
{Test} & \multicolumn{4}{c}{Weight $W$} \\
\cline{2-5}
{Case} & {1.0} & {1.2} & {2.0} & {5.0} \\
\midrule
0 & 205614 & 13113 & 3679 & 285 \\
1 & 5363 & 2243 & 828 & 288 \\
2 & 32541 & 4771 & 664 & 154 \\
3 & 50985 & 8587 & 1196 & 278 \\
4 & 707 & 358 & 1727 & 1801 \\
5 & 15099 & 7445 & 3841 & 235 \\
6 & 174411 & 54713 & 2060 & 276 \\
7 & 127135 & 31576 & 1189 & 199 \\
8 & 463 & 533 & 220 & 463 \\
9 & 1128 & 651 & 1063 & 250 \\
\midrule
Avg & 61344.6 & 12379.0 & 1646.7 & 422.9 \\
\bottomrule
\end{tabular}
\end{table}

\section{Analysis and Conclusion}

The results clearly demonstrate the trade-off between computational efficiency
and solution optimality inherent in Weighted A*:

\begin{itemize}
  \item \textbf{$W = 1.0$}: Produces the most optimal solutions (lowest cost)
  in all cases. However, the number of nodes explored is extremely high --
  sometimes exceeding 200,000 nodes for a single puzzle. This makes it
  impractical for time-sensitive applications or larger puzzle sizes.

  \item \textbf{$W = 5.0$}: Runs extremely fast, exploring only a few hundred
  nodes on average. However, the solution quality degrades significantly --
  costs are on average 60\% higher than optimal. The greedy nature of the
  search causes it to make locally optimal choices that lead to suboptimal
  global paths.

  \item \textbf{$W = 2.0$}: Offers a reasonable middle ground, with costs only
  23\% above optimal while reducing explored nodes by over 97\% compared to
  $W = 1.0$.

  \item \textbf{$W = 1.2$}: Strikes an excellent balance. The average cost is
  only 2\% above optimal, while the number of nodes explored drops by roughly
  80\% compared to $W = 1.0$. This makes $W = 1.2$ a strong candidate for
  practical use where near-optimal solutions are acceptable and computation
  time matters.
\end{itemize}

In conclusion, while $W = 1.0$ guarantees optimality, the computational cost
is prohibitive for complex puzzles. A small weight like $W = 1.2$ preserves
near-optimal solution quality while providing substantial speedups, making it
the recommended choice for most scenarios.

\end{document}
